% ══════════════════════════════════════════════════════════════════
% CONCLUSION GÉNÉRALE ET PERSPECTIVES
% ══════════════════════════════════════════════════════════════════
\newpage
\chapter*{Conclusion Générale et Perspectives}
\addcontentsline{toc}{chapter}{Conclusion Générale et Perspectives}

Ce stage de projet de fin d'études pendant 6 mois était intéressant et bénéfique pour moi, car j'ai pu mettre en pratique les connaissances théoriques et techniques acquises à l'université EMSI, en travaillant sur un projet concret et en contribuant activement à l'amélioration d'un système logiciel fourni par la société d'accueil.

Je suis à la fois fier et satisfait de mon stage au sein de la Société HPS Group. Car, ce stage m'a permis d'apprécier de plus près la vie professionnelle au sein d'une société, travailler en parfaite harmonie au sein d'une équipe et découvrir dans le détail comment gérer et réaliser un projet du début jusqu'à la fin.

Au sein d'HPS, j'ai pu m'ouvrir sur le domaine de test qui m'a permis de travailler sur des sujets qui partagent le même contexte. Ce projet m'a permis d'approfondir mes connaissances dans le métier de test avec une formation sur l'ISTQB et les sujets d'automatisation.

J'ai beaucoup apprécié les missions qui m'ont été confiés. Car, d'abord, elles s'inscrivent dans le cadre de mon projet personnel, ensuite elles m'ont permis de découvrir les étapes du processus de test, de l'expression des besoins jusqu'à la fin.

Tout au long de ce stage, j'avais un suivi et un encadrement de qualité, trois réunions par semaine avec l'équipe de test pour suivre mon avancement, et pour faire le point sur les points de blocage et tâches à venir.

Les difficultés que j'ai vécu durant mon stage sont principalement liées à l'incomplétude des cahiers de spécification et l'environnement technique de travail. Travailler sur cette difficulté m'a donné la possibilité de communiquer avec plusieurs personnes et de ne pas rester renfermée sur moi-même, et donc de s'intégrer plus facilement au sein de l'entreprise.

Pour conclure, ce stage de fin d'études a constitué une opportunité exceptionnelle d'évoluer au sein d'un environnement professionnel exigeant, centré sur l'innovation et la qualité des systèmes de paiement. La réalisation du Cypress Generator Dashboard et l'automatisation des tests fonctionnels ont répondu concrètement à un besoin d'optimisation des processus QA. L'intégration de l'intelligence artificielle (Mistral AI) pour la découverte automatique d'écrans et la génération de tests, combinée à l'orchestration des workflows via n8n, positionne cette solution comme un outil de nouvelle génération dans le domaine de l'assurance qualité logicielle.

Les connaissances acquises en ingénierie logicielle, couplées à une immersion dans les méthodologies de travail collaboratives, consolident ma formation d'ingénieur. Ce projet confirme l'importance stratégique de l'automatisation dans le cycle de vie logiciel moderne, ouvrant des perspectives prometteuses pour ma carrière future dans le domaine de la qualité et du développement logiciel.

\textbf{Perspectives d'amélioration :}
\begin{itemize}[leftmargin=*]
\item Intégration d'un modèle d'IA local (type LLM auto-hébergé) pour éliminer la dépendance aux API externes,
\item Extension du crawling AI à l'ensemble des modules PowerCARD pour couvrir 100\% des écrans,
\item Mise en place d'un pipeline CI/CD complet intégrant le Dashboard comme étape de génération automatique des tests de régression,
\item Développement d'un mode batch pour l'exécution parallèle de multiples scénarios de test,
\item Ajout d'un système de reporting avancé avec historique des exécutions et tendances de qualité.
\end{itemize}
